\subsection{Trees}

When one speaks of paths, or of connected components, of a graph, one means the paths or the connected components of the underlying quiver. Two vertices of a graph belong to the same connected component if and only if there exists a path connecting them.

Let G be a graph. If \(c = (a_{0}, f_{1}, a_{1}, \ldots, f_{n}, a_{n})\) is a path in G, the sequence \(\overline{c} = (a_{n}, \overline{f_{n}}, \ldots, a_{1}, \overline{f_{1}}, a_{0})\) is a path in G, called the opposite path to c. Let c and \(c'\) be juxtaposable paths in G; then, \(\overline{c'}\) and \(\overline{c}\) are juxtaposable, and one has \(\overline{c * c'} = \overline{c'} * \overline{c}\) .

A path c in G is said to be without backtracking if there is no pair of consecutive arrows of c that are opposite. Let \(c = (a_{0}, f_{1}, \ldots, f_{n}, a_{n})\) be a path in G connecting \(a_{0}\) and \(a_{n}\) . If, for an integer i such that \(1 \leqslant i < n\) , the arrows \(f_{i}\) and \(f_{i+1}\) are opposite, the path \((a_{0}, f_{1}, \ldots, a_{i-1}, f_{i+2}, \ldots, f_{n}, a_{n})\) is a path in G connecting \(a_{0}\) to \(a_{n}\) whose length is strictly less than that of the path c. By induction, there therefore exists a path without backtracking in G that connects \(a_{0}\) to \(a_{n}\) .

DEFINITION 2. --- A forest is a graph in which every loop without backtracking is a constant loop. A tree is a connected forest.

PROPOSITION 1. --- Let G be a graph. Every forest of G is contained in a maximal forest of G; in particular, there exists a maximal forest of G.

For a forest of G to be maximal, it is necessary and sufficient that its vertex set be equal to the vertex set of G and that its connected components be those of G.

Let \(A_{0}\) be a forest of G. The set of forests of G, equipped with the order relation “A is a subgraph of B,” is inductive. Consequently, there exists a maximal forest of G of which \(A_{0}\) is a subgraph (E, III, p. 21, cor. 1).

The subgraph of G whose vertex set is \(\operatorname{Som}(G)\) and whose arrow set is empty is a forest of G. Hence there exists a maximal forest of G.

Let A be a maximal forest of G. We prove that A and G have the same vertex set. The subgraph of G whose vertex set is \(\operatorname{Som}(G)\) and whose set of arrows is \(\operatorname{Fl}(A)\) is a forest and A is a subgraph of it. Thus we have \(\operatorname{Som}(A) = \operatorname{Som}(G)\) .

We now prove that A and G have the same connected components. Since every arrow of A is an arrow of G, each connected component of A is contained in a connected component of G. Since A and G have the same vertex set, it suffices to show that two vertices of G lying in the same connected component of G lie in the same connected component of A. If this is not the case, the relation \(R_{A}\) , "being in the same connected component of A," is strictly finer than the relation \(R_{G}\) and there exist two vertices of G that are not in the same connected component of A but are nevertheless joined by an arrow f of G. Let B be the directed subgraph of G whose vertex set is \(\operatorname{Som}(G)\) and whose arrow set is \(\mathrm{Fl(A)} \cup \{f, \overline{f}\}\) . Let us prove that B is a forest of G. Let \(c = (a_{0}, f_{1}, \ldots, f_{n}, a_{n})\) be a nonconstant loop without backtracking in B of minimal length. Since A is a forest, the loop c is not a loop in A. Let i (resp. j) be the smallest (resp. the largest) integer of \(\{0, \ldots, n\}\) such that \(a_{0}\) and \(a_{i}\) (resp. \(a_{j}\) and \(a_{n}\) ) are not in the same connected component of A. This means that the arrows \(f_{i}\) and \(f_{j+1}\) are oriented arrows of B associated with the edge \(\{f, \overline{f}\}\) and that they are opposite. Since the loop c has no back-and-forth, \(f_{i+1} \neq \overline{f_{i}}\) , therefore \(i \neq j\) and the path \((a_{i}, f_{i+1}, a_{i+1}, \ldots, f_{j}, a_{j})\) is a non-constant loop without backtracking in B of length < n, contrary to the hypothesis that c is of minimal length. It follows that B is a forest. This contradicts the hypothesis that A is a maximal forest of G.

Let now A be a forest of G such that \(\mathrm{Som}(\mathrm{A}) = \mathrm{Som}(\mathrm{G})\) and \(\pi_{0}(\mathrm{A}) = \pi_{0}(\mathrm{G})\) . Let us prove that this is a maximal forest of G. It suffices to show that if \(f \notin \mathrm{Fl}(\mathrm{A})\) , the subgraph B of G with vertex set \(\mathrm{Som}(\mathrm{G})\) and arrow set \(\mathrm{Fl}(\mathrm{A}) \cup \{f, \overline{f}\}\) is not a forest. By hypothesis, the points \(o(f)\) and \(t(f)\) lie in the same connected component of A; there therefore exists a path c without backtracking in A joining \(o(f)\) to \(t(f)\) . The path \(c * \overline{f}\) is then a loop without backtracking and non-constant in B, which shows that B is not a forest.

COROLLARY. --- A maximal forest of a connected graph is a maximal tree thereof.

Remark 1. --- In LIE, IV, p. 33, appendix, the notion of a combinatorial graph was defined as a pair (A, S), where S is a set and

A a subset of \(\mathfrak{P}(S)\) consisting of two-element sets; the elements of S are called vertices, those of A edges, and two vertices \(x\) and \(y \in S\) are said to be linked if \(\{x, y\}\) is an edge.

To such a combinatorial graph \(\Gamma = (\mathrm{A},\mathrm{S})\) , one associates a graph G whose vertex set is S and whose arrow set \(\widetilde{\mathrm{A}}\) is the subset of \(\mathrm{S}^2\) formed by the pairs of linked vertices, with the source map and the target map coinciding with the first and second projections of \(\mathrm{S}^2\) in S, and the involution \(f\mapsto \overline{f}\) being given by the restriction to \(\widetilde{\mathrm{A}}\) of the map \((x,y)\mapsto (y,x)\) of \(\mathrm{S}^2\) into itself. The map which to an edge \(\{f,\overline{f}\}\) of G associates the set \(\{o(f),t(f)\}\) is a bijection from the set of edges of G onto that of the edges of the combinatorial graph \(\Gamma\) .

Conversely, every graph such that the source and target of every arrow are distinct, and such that an arrow is determined by its source and target, is of this form.

The reader will verify that the notions of connectedness, tree, or forest for a combinatorial graph coincide with the corresponding notions for the graph associated with it.

